% !TeX root = théorie.tex

\section{Semi-groupe}

\begin{definition}
Soit $E$ un ensemble et $*: E \times E \rightarrow E$ un loi interne.
$(E, *)$ est un semi-groupe si $*$ est associative.
$(E, *)$ est un monoïde si de plus, il existe un élément neutre.
% unifère
\end{definition}

\begin{definition}[Idéal]
Soit $(E, *)$ un semi-groupe.
L'opération $*$ peut être étendue aux parties de $E$ ainsi,
pour $(A, B) \in \mathcal{P}(E)^2$ :
\begin{equation}
    A * B = \left\{ a * b \;|\; (a, b) \in A \times B \right\}
\end{equation}
Soit $I \subset E$. On dit que $I$ est
un idéal à droite si $I * E \subset I$
et un idéal à gauche si $E * I \subset I$.
Si $(E, *)$ est commutatif, on dit simplement idéal puisque $I*E = E*I$.
\end{definition}

\begin{proposition}
Soit $(E, *)$ un semi-groupe fini.
Il existe un élément de $E$ idempotent.
\end{proposition}

\begin{proof}
À faire
\end{proof}

\begin{theoreme}
Soit $(E, *)$ un semi-groupe commutatif.
Si $I$ l'idéal minimal de $E$ (l'intersection de tous les idéaux)
alors $(I, *)$ est un groupe commutatif.
\end{theoreme}
